\documentclass[10pt,a4paper]{amsart}
\usepackage[margin=19mm]{geometry}
\usepackage{amsmath,amssymb,mathtools}
\usepackage[scr=rsfs,cal=euler]{mathalfa}
\usepackage{xfrac}
\usepackage[unicode,hidelinks,bookmarks=false]{hyperref}
\newcommand{\ie}{i.e.\ }
\newcommand{\cf}{cf.\ }
\newcommand{\resp}{respectively\ }
\newcommand{\Subsection}{\subsection}
\providecommand{\nobreakdash}{\nobreak\hbox{-}}
\setlength{\emergencystretch}{2em}
\multlinegap0pt
\usepackage{amssymb}
\newtheorem{theorem}{Theorem}
\newtheorem{lemma}{Lemma}
\newcommand{\area}{\Delta}
\title{\textbf{Average Chord Lengths in a Triangle}}
\author{Stanley Rabinowitz}
\date{Source: arXiv:2609.11370v1 (CC BY 4.0)}
\begin{document}
\maketitle

\noindent\emph{Source and licence.} \textbf{Average Chord Lengths in a Triangle}. Version arXiv:2609.11370v1 (CC BY 4.0), distributed by arXiv under the Creative Commons Attribution 4.0 International licence, http://creativecommons.org/licenses/by/4.0/, as recorded in the arXiv licence metadata for this exact version and shown on the abstract page https://arxiv.org/abs/2609.11370v1. Full licence notice: \texttt{licenses/arxiv-2609.11370-CC-BY-4.0.txt}.\par\smallskip

\noindent\textit{Editorial scope.} Counted unit: the article's theorem thm:main (source0.tex lines 400--411). The article's complete original proof of it is delivered with it (source0.tex lines 514--588, one \texttt{proof} environment of 75 lines, which the article places away from the statement). No mathematics is written, repaired or re-derived: the statement and the proof are the article's own lines, in the article's own notation, and no retained line is substituted or re-flowed.  Retained uncounted as the source-local context the proof needs: lines 416--433; lines 486--492.  Nothing was omitted from any retained span, so the package carries no omission record.  No retained line contains a citation, so the wrapper carries no bibliography and no bibliography entry.  No retained line contains a figure environment, an image-inclusion command or drawing code, so no image is delivered and the package declares no asset dependency.  Editorial formatting only, in addition: an AMS \texttt{amsart} preamble that restates the article's own declarations and macros used by the retained lines, namely the environments \texttt{theorem}, \texttt{lemma} on separate counters, together with the standard packages those lines need.  The retained environments print in this document as Theorem 1, Lemma 1 in order of appearance, whereas the article prints the same items with the numbering of its own full document.  The complete native source of the article as distributed in the arXiv e-print for this exact version is bundled unchanged as \texttt{sources/210012/source0.tex}.
\begin{lemma}\label{lem:symm}
Fix a side $BC$ of a triangle $T$ and a point $P\in T$.  Replace
each segment of $T$ parallel to $BC$ by a segment of the same length
whose midpoint lies on the line through $P$ perpendicular to $BC$.
The resulting set $T^*$ is an isosceles triangle with the same base
length and altitude as $T$, and $P\in T^*$.  Moreover,
\[
 \int_{T^*}\frac{dA(X)}{|X-P|}
 \ge
 \int_T\frac{dA(X)}{|X-P|},
\]
and
\[
 p(T^*)\le p(T).
\]
Equality in the first inequality holds only when $T$ is already
symmetric about the line through $P$ perpendicular to $BC$.
\end{lemma}

\begin{lemma}\label{lem:area-bound}
For every triangle $T$ of area $\area$ and every point $P\in T$,
\[
 \int_T\frac{dA(X)}{|X-P|}
 \le 2\sqrt{\pi\area}.
\]
\end{lemma}

\begin{theorem}\label{thm:main}
Let $T$ be a triangle of perimeter $p$, and let $P$ be any interior
point of $T$.  Then
\[
 {
 M_T(P)\le
 \frac{p}{\pi\sqrt3}\log(2+\sqrt3).
 }
\]
Equality holds if and only if $T$ is equilateral and $P$ is its
center.
\end{theorem}

\begin{proof}[Proof of Theorem~\ref{thm:main}]
Because both perimeter and $M_T(P)$ scale linearly, fix the perimeter
$p$.  It is convenient here to allow $P$ temporarily to lie on the
boundary of $T$; the integral
\[
 V(T,P)=\int_T\frac{dA(X)}{|X-P|}
\]
is still finite.

A maximizing pair $(T,P)$ exists.  Indeed, after translating $P$ to
the origin, all vertices lie in a fixed bounded disk.  A maximizing
sequence therefore has a convergent subsequence.  A degenerate
limiting triangle has area tending to zero, and
Lemma~\ref{lem:area-bound} then forces $V(T,P)\to0$, so a maximizer
cannot be degenerate.  For a nondegenerate limit, continuity follows
by translating $P$ to the origin: the indicator functions of the
triangles converge almost everywhere in a common bounded disk, while
$1/|X|$ is locally integrable in the plane.  Thus the maximum is
attained on a nondegenerate triangle.  The argument below will in
fact force the maximizing point to be the interior center of an
equilateral triangle.

Choose any side $BC$ and form the triangle $T^*$ of
Lemma~\ref{lem:symm}.  Then
\[
 V(T^*,P)\ge V(T,P),
 \qquad
 p(T^*)\le p.
\]
Set
\[
 \lambda=\frac{p}{p(T^*)}\ge1
\]
and enlarge $T^*$ by a homothety centered at $P$ with factor
$\lambda$.  Since
\[
 V(\lambda T^*,P)=\lambda V(T^*,P)
 \ge \lambda V(T,P),
\]
maximality forces $\lambda=1$ and
\[
 V(T^*,P)=V(T,P).
\]
Thus equality holds throughout Lemma~\ref{lem:symm}.  The perimeter
equality gives $AB=AC$, while equality in the integral shows that
$T$ is symmetric about the line through $P$ perpendicular to $BC$;
in particular, $P$ lies on the perpendicular bisector of $BC$.

Apply the same argument to a second side.  The triangle is then
isosceles with respect to two different sides, and hence equilateral.
The two symmetry axes meet at its center, so $P$ is that center.

It remains only to evaluate $M_T(P)$ for an equilateral triangle.
Let its side length be $a=p/3$.  The distance from its center to a
side is
\[
 d=\frac{a\sqrt3}{6}=\frac{p}{6\sqrt3}.
\]
For the rays meeting one side,
\[
 \rho(\theta)=d\sec\theta,
 \qquad -\frac{\pi}{3}\le\theta\le\frac{\pi}{3}.
\]
Thus
\[
 \int_0^{2\pi}\rho(\theta)\,d\theta
 =
 3d\int_{-\pi/3}^{\pi/3}\sec\theta\,d\theta
 =
 6d\log(2+\sqrt3)
 =
 \frac{p}{\sqrt3}\log(2+\sqrt3).
\]
Division by $\pi$ completes the proof.
\end{proof}

\end{document}
